\documentclass[11pt,a4paper]{article}
\usepackage{fontspec}
\setmainfont{Noto Serif}
\setsansfont{Noto Sans}
\usepackage{babel}
\babelprovide[main,import]{russian}
\usepackage{amsmath,amssymb,mathtools,geometry,fancyhdr,microtype,xcolor,hyperref}
\geometry{left=21mm,right=21mm,top=20mm,bottom=21mm}
\definecolor{ink}{HTML}{173D57}
\hypersetup{colorlinks=true,urlcolor=ink,linkcolor=ink,pdftitle={BGS-F. Перенос между торсионными секторами}}
\pagestyle{fancy}\fancyhf{}\lhead{\small BGS-F}\rhead{\small Перенос между торсионными секторами}\cfoot{\thepage}
\setlength{\headheight}{15pt}\setlength{\parskip}{2.3pt}\setlength{\emergencystretch}{3em}
\newcommand{\Tr}{\operatorname{Tr}}
\newcommand{\Ad}{\operatorname{Ad}}
\newcommand{\Op}{\operatorname{Op}}
\newcommand{\dist}{\operatorname{dist}}
\newcommand{\T}{\mathbb{T}}
\newcommand{\Z}{\mathbb Z}
\newcommand{\C}{\mathbb C}
\newcommand{\Qbar}{\overline{\mathbb Q}}
\newcommand{\sect}[1]{\section*{\sffamily\color{ink}#1}}
\newcommand{\subt}[1]{\par\medskip\noindent\textbf{#1.}\quad}
\begin{document}
\raggedright
\noindent{\small\sffamily\color{ink} МАТЕМАТИЧЕСКАЯ ФИЗИКА\hfill 22 НОЯБРЯ 2026}
\begin{center}{\LARGE\sffamily Перенос между торсионными секторами}\par\vspace{4pt}
{\large Полный фиксированный кик на точных уровнях Вейля}
\end{center}

Пусть \(N\) нечётно, \(L\) нечётно, \(L^2\mid N\), \(3\le L\le N^{1/20000}\). Совместные собственные пространства двух коммутирующих переводов Вейля на \(N/L\) обозначим через \(P_\chi\), где \(\chi\) пробегает \(L^2\) характеров. Каждый проектор имеет ранг \(d=N/L^2\). Для исходной матрицы
\[
 U_N=\exp(i\tau N A_N)M_N,\qquad \tau=10^{-80},
\]
справедливо равномерно по таким \(L,\chi,\psi\)
\begin{equation}
    \frac{1}{d}\Tr(P_\chi U_N P_\psi U_N^*)
     =\frac{1}{L^2}\big[1+O_\tau(N^{-1/6000})\big].                 \tag{1}
\end{equation}
То же верно для одного кика \(\exp(i\tau N A_N)\). Это средняя доля квадрата нормы из входного сектора, попавшая в заданный выходной сектор. Для канонического сектора при фиксированном \(L\) предел \(L^{-2}\) был известен. Здесь контролируются все пары характеров со степенной ошибкой. Например, на \(N=7^k\) можно взять \(L=7^{\lfloor k/20000\rfloor}\). Константа при фиксированном \(\tau\) не оптимизирована. Эффективный порог размера здесь не заявлен.

\subt{Рассеяние больших меток Вейля}
Для всех нечётных \(N\) и метки \((a,b)\), чья центрированная величина
\(\max(|a|,|b|)\ge N^{99/100}\), коэффициенты разложения \(U_N T_N(a,b)U_N^*\)
в базисе Вейля по модулю равномерно не больше \(C_\tau N^{-9/64000}\).
На выходную метку и делители \(N\) ограничений нет.

\subt{Точное отсутствие сохранённых векторов}
Для всяких нечётных \(N,L\) с \(L^2\mid N\), \(L\ge3\), и любого вещественного алгебраического \(t\ne0\) имеем
\begin{equation}
  \operatorname{Ran}(P_\psi)\;\cap\;
   \big(e^{it A_N}M_N\big)^*\operatorname{Ran}(P_\chi)=\{0\}.
                                                                    \tag{2}
\end{equation}
Аналогично для самого кика. Здесь ограничения \(L\le N^{1/20000}\) нет. Точное вложение нижнего метаплектического уровня через неподвижный характер перестаёт быть вложением для полного оператора.

\subt{Связь с приложенной запиской}
Секторальное предложение приводит к независимой оценке матрицы перекрытий двух квантовых разложений, с использованием именно заданных коэффициентов \(A_N\). Оценка получена при неизменной силе для полного экспоненциального оператора, при неограниченном росте \(\tau N\). В последнем разделе проверяется, что этих перекрытий недостаточно для парной рампы \((P)\). У исходного оператора остаются открытыми и \((P)\), и локальный предел \(\mathrm{Sine}_2\).

\newpage
\sect{1. Точные уровни и квантовый символ}
Положим \(e_N(s)=e^{2\pi i s/N}\) и
\[
 (T_N(a,b)v)(x)=e_N\big(b(x+\overline2 a)\big)v(x+a).
\]
Типы квазиимпульса и положения в этом определении совпадают с постановкой BGS-F. Для \(z=(a,b)\in(\Z/L)^2\) зададим \(W_z=T_N((N/L)a,(N/L)b)\). В фазе произведения \(W_zW_{z'}\) появляется множитель \(e_N(\overline2(N/L)^2(ab'-ba'))\), равный 1 из-за \(L^2\mid N\). Это представление группы \((\Z/L)^2\), с \(\Tr W_z=N\) при \(z=0\) и нулём иначе. Поэтому
\begin{equation}
   P_\chi=\frac1{L^2}\sum_z\overline{\chi(z)}W_z,\qquad
   \Tr P_\chi=d,\qquad P_\chi P_\psi=\mathbf1_{\chi=\psi}P_\chi.
                                                                  \tag{3}
\end{equation}
Точное метаплектическое сопряжение переводит \(W_z\) в \(W_{Sz}\) (при обратном порядке сопряжения в \(W_{S^{-1}z}\)). Значит, \(M_N\) переставляет проекторы. При тривиальном характере \(P_1\) образ состоит из функций с носителем на \(L\Z/N\Z\), постоянных при сдвиге на \(N/L\). Изометрия \(J\) из уровня \(d\) сопоставляет базисному вектору \(y\) сумму векторов в положениях \(Ly+jN/L\), \(0\le j<L\), с нормировкой \(L^{-1/2}\). Точное равенство \(M_NJ=JM_d\) видно из ядра \(M_N\). Зависимость от \(j\) в фазе сводится к \(e_L(2\overline8 xj)\). Сумма исчезает при \(L\nmid x\), а при \(x=Lx'\) оставшаяся фаза равна фазе \((M_d)_{x',y}\). Нормировка \(L^{1/2}/\sqrt N=L^{-1/2}/\sqrt d\) завершает проверку.

Для \(f\in C^\infty(\T^2)\) используем обычную полуклассическую конвенцию
\[
 \Op_N(f)=\sum_{k,l\in\Z}\widehat f(k,l)(-1)^{kl}T_N(l,k).
\]
Множитель \((-1)^{kl}\) переводит модульную \(\overline2\) в арифметическую половину в фазе. Так, \(A_N=\Op_N(g)\) с
\begin{align*}
 g(q,p)&=\cos(14\pi q)-1/10+
 (1-\cos(14\pi q))(1+\cos(2\pi q))
       (1/10+\sin(2\pi q)/100)\\
 &\quad+\tfrac1{100}r(q+1/2)\cos(2\pi p-\alpha_0),\\
 r(q)&=(1-\cos(14\pi q))(1+\cos(2\pi q)).
\end{align*}
Коэффициенты \(\widehat g(1,0)=1/20-i/200\) и
\(\widehat g(0,1)=e^{-i\alpha_0}/200\) оба ненулевые. Носитель удовлетворяет \(|k|\le9, |l|\le1\), и \(\sum |\widehat g(k,l)|<2\), по треугольному неравенству из разложения выше. Под номером (1) будет использовано, что \(W_z\) квантует торический сдвиг \(h_z=(-a/L,b/L)\). Это следует прямо из сопряжения \(W_z^*\Op_N(f)W_z=\Op_N(f(\cdot+h_z))\).

\newpage
\sect{2. Количественный след через функцию смещения}
Обозначим \(\operatorname{tr}_N=\Tr/N\) и \(\|X\|_{2,N}^2=\operatorname{tr}_N(X^*X)\). Пусть \(K_N=e^{i\tau N A_N}\) и \(\varphi\) есть классический поток кика за физическое время \(\tau\). Он аналитичен и сохраняет площадь. Его поле \(Y\) имеет компоненты \((-g_p,g_q)/(2\pi)\), в соответствии с операторной конвенцией \(K_N^*\Op_N(f)K_N\). В частности модули ненулевых коэффициентов \(Y\) на двух указанных частотах равны модулям соответствующих коэффициентов \(g\) в одной из компонент.

\subt{Контроль символов}
Для гладких функций \(a,b\) фазовое тождество Вейля и \(|e^{ix}-1|\le |x|\) дают
\[
 \|\Op_N(a)\Op_N(b)-\Op_N(ab)\|
       \le C N^{-1}\|a\|_{C^6}\|b\|_{C^6}.
\]
Действительно, каждая ошибка в свёртке ограничена
\(C|\nu||\mu|/N\), где \(\nu,\mu\in\Z^2\) — частоты, и взвешенный ряд Фурье для нормы \(C^6\) суммируется. В коммутаторе с \(A_N\) линейный член синуса даёт \(Yf\) для производной \(-iN[A_N,\Op_N(f)]\). Остаток оценивается \(|\sin x-x|\le |x|^3/6\) и суммой \(\sum_\nu |\nu|^3|\widehat f(\nu)|\), ограниченной \(C\|f\|_{C^8}\). Интегрирование в операторной формуле Дюамеля вдоль гладкого потока даёт
\[
 \|K_N^*\Op_N(f)K_N-\Op_N(f\circ\varphi)\|
     \le C_\tau N^{-2}\|f\|_{C^8}.
\]
Эта же оценка верна для обратного кика с обратным потоком. Наконец, \(\operatorname{tr}_N\Op_N(a)=\int a+O(N^{-4}\|a\|_{C^4})\): ненулевые следы отдельных частот возможны лишь для \(\nu\in N\Z^2\). С оценкой произведения это даёт
\[
  \|\Op_N(a)\|_{2,N}\le\|a\|_{L^2(\T^2)}
                              +C N^{-1/2}\|a\|_{C^6}.             \tag{4}
\]

\subt{Усреднение малыми переводами}
Возьмём \(V=W_z K_N W_{z'}K_N^*\), \(G=T_{h_z}\varphi T_{h_{z'}}\varphi^{-1}\), где \(T_h\) здесь означает сдвиг точки тора. Положим \(B_\nu=\Op_N(e^{2\pi i\nu\cdot x})\), унитарный оператор. Для \(R\ge1\) усредним \(B_\nu^* V B_\nu\) по квадрату \(|\nu|_\infty\le R\), получая \(J_R\). \(\operatorname{tr}_N J_R=\operatorname{tr}_N V\). После домножения слева на \(V^*\) получаем через Егорова и оценку произведения
\begin{align*}
 \|V^*J_R-\Op_N(a_R)\|&\le C_\tau N^{-1}(1+R)^8,\\
 a_R(x)&=D_R(x_1-G_1(x))D_R(x_2-G_2(x)),\\
 D_R(u)&=\frac1{2R+1}\sum_{j=-R}^R e^{2\pi i j u}.
\end{align*}
В оценке произведения один из сомножителей есть отдельная вейлева гармоника, поэтому вместо её \(C^6\)-нормы достаточно множителя \(C(1+R)\) из фазовой формулы. Другой множитель контролируется в \(C^6\) и стоит \(O_\tau((1+R)^6)\). Все производные \(G\) требуемых порядков ограничены независимо от сдвигов. В частности \(\|a_R\|_{C^6}\le C_\tau(1+R)^6\). По (4) и неравенству Коши
\begin{equation}
 |\operatorname{tr}_N V|\le \|a_R\|_2+
        C_\tau\big(N^{-1/2}(1+R)^6+N^{-1}(1+R)^8\big).            \tag{5}
\end{equation}
Использована квантовая фаза полного оператора. Здесь нет оценки через интенсивности классической цепочки или через неизвестные следы \(U_N^m\).

\newpage
\sect{3. Равномерный контроль малых смещений}
Запишем \(h=h_z\), \(u=h_{z'}\), центрируя обе координаты в \([-1/2,1/2]\). При \(u\ne0\) подставим \(y=\varphi(x)\) в \(G(y)-y\). В поднятии на плоскость
\begin{equation}
    G(\varphi(x))-\varphi(x)=
   h+u+\tau[Y(x+u)-Y(x)]+E(x+u)-E(x),\qquad
   E=\varphi-\mathrm{id}-\tau Y.                                  \tag{6}
\end{equation}
Подстановка сохраняет площадь. Если \(j\) максимизирует \(|u_j|\), берём частоту \(\nu=e_j\) и ту компоненту поля, в которой соответствующий коэффициент имеет модуль \(\ge1/200\). Из \(|e^{2\pi i u_j}-1|\ge4|u_j|\) амплитуда выбранного коэффициента при первом порядке не меньше \(\tau |u|_\infty/50\).

Для контроля полной ошибки используем полосу \(|\operatorname{Im}x_i|<1/8\), где из носителя ряда \(g\) следует \(\|Y\|_\infty<6\cdot10^4\), \(\|DY\|_\infty<4\cdot10^6\) в максимальной норме. Комплексный поток из полосы \(1/16\) остаётся в полосе \(1/8\) на \([0,\tau]\) по интегральному уравнению. Отсюда \(\|E\|_{1/16}\le10^{12}\tau^2\), так как \(\varphi^s-x\) не превосходит \(6\cdot10^4 s\). Индекс нормы обозначает ширину полосы. По формуле Коши и интегрированию производной по отрезку \([x,x+u]\)
\[
                   \|E(\cdot+u)-E\|_{1/20}
                          \le10^{15}\tau^2 |u|_\infty .          \tag{7}
\]

Пусть \(f(x)\) есть выбранная компонента в (6). Обозначим \(D_q=(2\pi i)^{-1}\partial_q\) и аналогично \(D_p\). Дифференциальный оператор Лагранжа
\[
  \mathcal L_\nu=
  \prod_{a=-9,\ a\ne\nu_1}^9\frac{D_q-a}{\nu_1-a}
  \prod_{b=-1,\ b\ne\nu_2}^1\frac{D_p-b}{\nu_2-b}
\]
выбирает ровно одну частоту из прямоугольника носителя, подавляя и константу. Порядок равен 20, сумма модулей коэффициентов ограничена \(((10!)^2)4<10^{14}\), поскольку знаменатели — ненулевые целые. Для производных \(D_q^aD_p^b\) порядка до 20 на вещественном торе оценка Коши в полидиске радиуса \(1/32\) даёт множитель не больше \(20!(32/(2\pi))^{20}<10^{33}\). Поэтому эта фильтрация ошибки (7) стоит не более \(10^{62}\tau^2|u|_\infty\). При \(\tau=10^{-80}\) получаем в каждой точке
\[
          |\mathcal L_\nu f(x)|\ge\tfrac1{100}\tau|u|_\infty .
\]
У оператора \(\mathcal L_\nu\) нулевой свободный коэффициент. Значит, некоторая действительная производная положительного порядка \(\le20\) по модулю не меньше \(c\tau|u|_\infty\) в данной точке. Все положительные производные \(f\) до 21-го порядка по (6), (7) ограничены \(C\tau|u|_\infty\) равномерно по обоим сдвигам.

\newpage
\sect{4. Оценка перекрытий при растущем числе уровней}
\subt{Оценка меры}
Из производной порядка \(a\) в предыдущем разделе следует нижняя граница \(c'\tau|u|_\infty\) для \(|\partial_v^a f|\) в одном из конечного набора направлений. Действительно, однородный многочлен производных можно восстановить по значениям на \((1,j/20)\), \(0\le j\le20\), по одномерной интерполяции. Граница на порядок \(a+1\) позволяет выбрать радиус окрестности независимо от \(h,u\). Фиксированная конечная сетка центров покрывает тор такими окрестностями. На каждой используем сечение в направлении \(v\). Если на отрезке \(|\partial_v^a f|\ge d_0\), множество \(|f-j_0|\le\varepsilon\) имеет длину \(\le C_a(\varepsilon/d_0)^{1/a}\). Чтобы это увидеть, из множества меры \(\ell\) выбирают \(a+1\) точек с расстояниями между соседями \(\ge\ell/(2a)\) при помощи квантилей. Формула разделённых разностей и теорема о среднем ограничивают \(d_0/a!\) сверху величиной \((a+1)\varepsilon(2a/\ell)^a\). Возможных целых \(j_0\) конечное равномерное число по (6). После интегрирования сечений получаем
\[
   |\{x:\dist(f(x),\Z)\le\varepsilon\}|
         \le C\left(\frac{\varepsilon}{\tau |u|_\infty}\right)^{1/20}
         \qquad (0<\varepsilon\le1).
\]
Если отношение больше 1, требуемое неравенство следует из полной площади 1. По геометрической сумме \(|D_R(s)|\le\min(1,C/(R\dist(s,\Z)))\). В (5) используем границу меры при \(\varepsilon=R^{-1/2}\) и \(|D_R|\le1\) для второй координаты. Это даёт
\[
  \|a_R\|_2\le C(\tau|u|_\infty)^{-1/40}R^{-1/80}+CR^{-1/2}.
\]
При \(R=\lfloor N^{1/32}\rfloor\) получаем для всякого \(z'\ne0\)
\begin{equation}
   |\operatorname{tr}_N(W_z K_N W_{z'}K_N^*)|
                 \le C_\tau L^{1/40}N^{-1/2560}.                  \tag{8}
\end{equation}
Для \(z'=0,z\ne0\) след в (8) равен нулю точно. Аналогично он равен нулю при \(z=0,z'\ne0\) по цикличности. Константы в (8) не требуют считать \(L\) фиксированным.

Тот же вывод до подстановки \(|u|_\infty\ge1/L\) даёт \(C_\tau |u|_\infty^{-1/40}N^{-1/2560}\) для произвольных квантованных сдвигов Вейля. При входном \(|u|_\infty\ge N^{-1/100}\) имеем \(1/2560-1/4000=9/64000\). Действие \(M_N\) меняет расстояние центрированного сдвига до нуля не более чем в фиксированное число раз, поскольку \(S,S^{-1}\) целочисленны. Тем самым доказана оценка для всех больших меток \(U_N\).

Подставляя (8) в два разложения (3), получаем
\[
 \operatorname{tr}_N(P_\chi K_N P_\psi K_N^*)
       =L^{-4}+O_\tau(L^{1/40}N^{-1/2560}).
\]
Для относительной оценки достаточно \(L^{4+1/40}N^{-1/2560}\le N^{-1/6000}\), что выполняется при \(L\le N^{1/20000}\) по сравнению рациональных показателей. Это доказывает (1). В ортонормированном по \(\operatorname{tr}_N\) базисе \(W_z\), \(z\ne0\), сжатие \(\Ad_{U_N}\) на бесследовую часть этой алгебры имеет операторную норму \(\le C_\tau L^{2+1/40}N^{-1/2560}\). Это следует из (8) заменой \(z\) на \(-z\), перестановкой столбцов от \(M_N\) и границей суммы модулей в строках и столбцах.

\newpage
\sect{5. Точный ранг утечки, включая фиксированное \(\tau N\)}
В базисе положений \(A_N\) — эрмитова циклическая матрица Якоби. На ребре \(x,x+1\) модуль коэффициента равен
\[
           \frac1{200}\,r\!\left(\frac{x+1/2}{N}+\frac12\right).
\]
Формула непосредственно следует из двух гармоник \(l=\pm1\) и перевода \(\overline2\) в обычную половину. Нули в физической координате середины ребра — \(0\) и \(1/2+j/7\) по модулю единицы. Первый недостижим на сетке \((x+1/2)/N\). Остальные разнесены на \(1/7\), поэтому оба ребра у одного положения одновременно не могут иметь нулевой коэффициент при \(N\ge9\).

Сектор \(P_\psi\) имеет носитель на одном вычете \(x\bmod L\), заданном характером модуляции \(W_{(0,1)}\). Поскольку \(L\ge3\), два соседних вычета различны между собой и отличны от исходного. Для любой \(v\) из входного сектора две внедиагональные части \(A_Nv\) поддержаны на этих соседних вычетах. Отдельное значение \(v(x)\) входит в соответствующую компоненту каждого соседа без суммирования с другими входными значениями. У каждого \(x\) хотя бы один коэффициент ненулевой. Поэтому \((I-P_\psi) A_NP_\psi\) инъективен на \(\operatorname{Ran}P_\psi\).

Рассмотрим \(B(t)=(I-P_\chi)e^{itA_N} P_\psi\) как прямоугольное отображение из входного сектора. При \(\chi\ne\psi\) оно инъективно при \(t=0\), при \(\chi=\psi\) инъективна производная в нуле, а \(B(0)=0\). В обоих случаях какой-то минор размера \(d\) в алгебраическом базисе является нетождественно нулевой аналитической функцией. Проекторы имеют алгебраические элементы. По спектральному разложению эрмитовой алгебраической матрицы \(A_N\) минор есть экспоненциальный многочлен
\[
                         \sum_\lambda c_\lambda e^{it\lambda},
           \qquad c_\lambda,\lambda\in\Qbar .
\]
После группировки частот остаётся хотя бы один ненулевой коэффициент. При вещественном алгебраическом \(t\ne0\) значения \(it\lambda\) — попарно различные алгебраические числа. Линейная независимость их экспонент над \(\Qbar\) по теореме Линдемана--Вейерштрасса даёт ненулевой минор в точном \(t\). Отсюда утверждение для кика. Перестановка проекторов матрицей \(M_N\) переносит вывод на (2). Доказательство минором не даёт равномерной нижней оценки всех сингулярных чисел.

\subt{Интерпретация}
На точных уровнях \(N=L^2d\) поднятие линейной карты существует как обычное вложение. Полный кик разрушает его в двух количественно различных смыслах: ни один ненулевой вектор не остаётся целиком внутри заданного сектора при алгебраическом времени, а средняя масса проекции в каждый сектор подчиняется (1). Второе утверждение использует масштаб \(t=\tau N\) и не вытекает из первого.

\subt{Предыдущие результаты}
Для канонического сектора точная утечка при алгебраическом времени и предел средней сохранённой массы \(L^{-2}\) при фиксированном \(L\) доказаны в
\href{https://github.com/ChurilovMaksim/BGS/blob/e2757e4cc5ef996e08190f88c93c710377c10c22/mathematics/spectral_research/proofs/research_round2_2026_10_20/arithmetic_oldspace_escape.md}{записке от 20 октября, теоремы 1 и 3}.
В
\href{https://github.com/ChurilovMaksim/BGS/blob/e2757e4cc5ef996e08190f88c93c710377c10c22/mathematics/fixed-map/seven_oldspace_singular_spectrum.md}{заметке от 18 ноября}
для \(L=7\) получена сходимость множества сингулярных чисел канонического сжатия к \([0,1]\). Формулы (1) и (2) охватывают разные пары характеров. В (1) контролируются растущие \(L\) и степенная ошибка. Оценка (2) исключает сингулярное число ровно 1 при каждом конечном \(N\), что совместимо с пределом нормы 1.

\newpage
\sect{6. Проверка рамочной и резольвентной постановок}
Для обычного произведения Гильберта--Шмидта \(\langle X,Y\rangle=\Tr(X^*Y)\) обозначим \(v_m=U_N^m/\sqrt N\), \(e=v_0\) и \(Q_w=\sum_m w(m/N)|v_m\rangle\langle v_m|\). Тогда
\begin{align*}
 \langle e,Q_w e\rangle
   &=\sum_m w(m/N)|\langle e,v_m\rangle|^2
     =N^{-2}\sum_m w(m/N)|\Tr U_N^m|^2,\\
 \Tr_{\mathrm{HS}}Q_w^2
   &=\sum_{m,n}w(m/N)w(n/N)|\langle v_m,v_n\rangle|^2\\
   &=N^{-2}\sum_{m,n}w(m/N)w(n/N)|\Tr U_N^{n-m}|^2 .
\end{align*}
Так проверяются (F1), (F3) записки. При фиксированном \(m\) полные \(T_N(a,b)U_N^m/\sqrt N\) образуют ортонормированный базис, поскольку умножение справа на \(U_N^m\) — изометрия. Сумма их проекторов точно равна \(I_{\rm HS}\) для всякого \(U_N\), в частности для clock и для CUE. Это (NG).

Для \(D=\tfrac{N^2}{4}(2I-\Ad_{U_N}-\Ad_{U_N}^*)\) положим \(R_s=(s^2+D)^{-1}\), \(s>0\), и \(P\) — проектор в HS на алгебру из (3), \(Q=I-P\). Точная формула Шура даёт
\[
   P R_s P =
   [s^2+PDP-PDQ(s^2+QDQ)^{-1}QDP]^{-1}
   \quad\text{на }\operatorname{Ran}P .
\]
Вычитаемая самоэнергия положительна, а обратная положительная матрица не больше \(s^{-2}I\). Поэтому вклад этой части в \(N^{-1}\Tr_{\mathrm{HS}}R_s\) ограничен \(L^2/(Ns^2)\), что стремится к нулю на доказанном диапазоне \(L\). Полная резольвента имеет вклад как минимум \(1/s^2\) из диагональных собственных операторов. След дополнения \(Q R_s Q\) требует отдельной оценки. Положительность самоэнергии в малом секторе в частности не даёт коэффициента отрицательной поправки \((P')\).

\subt{Контрпример для одних секторных перекрытий}
Отождествим пространство, сохраняя секторы, с \(\C^{L^2}\otimes\C^d\) и возьмём \(Z=F_{L^2}\otimes I_d\), где \(F_{L^2}\) — нормированное дискретное преобразование Фурье. Для него (1) выполняется без ошибки и бесследовое сжатие \(\Ad_Z\) на ту же алгебру равно нулю точно. Однако \(Z^4=I\), и ненулевая часть парной суммы при \(4\mid m\) растёт линейно по \(N\) для исходного \(w\), положительного на внутреннем интервале. При этом вычитаемый CUE-эталон равномерно ограничен. Следовательно, даже равномерность всех одношаговых перекрытий совместима с неправильными гейзенберговскими следами.

\subt{Воспроизведение и объём}
\texttt{torsion\_transport/check.py} проверяет точные числа для символа, константы и показатели степеней, представление секторов и формулу Якоби. Сценарий включает численный контроль формулы для фиксированной существенно большей тестовой силы, чтобы ненулевой эффект не исчезал в округлении. Авторская аналитическая оценка для \(10^{-80}\) приведена в разделах 2--4. Исходная записка сохранена в архиве для сопоставления.
\end{document}
